% GENERATED FILE. Edit canonical lessons, then run npm run generate:books.
% canonical-source-hash: fnv1a64:921c9bddb31b94be
% canonical-lessons: JA-C400-tokonoma, JA-C400-zashiki, JA-C400-engawa, JA-C400-otearai, JA-C400-shimeru

\chapter{A Japanese-Style Room}
\label{ch:ja-a-japanese-style-room}
\hlchaptermodality{400}

\begin{chapteropening}
\textbf{By the end of this chapter:} \emph{I can describe a traditional Japanese room.}

Complete the last lesson of chapter 400: I can describe a traditional Japanese room.
\end{chapteropening}

\section[\texorpdfstring{\ja{とこのま} (tokonoma)}{とこのま (tokonoma)}]{\ja{とこのま} (tokonoma) — an alcove (for flowers or a hanging scroll)}

\label{lesson:JA-C400-tokonoma}



\begin{quote}
Before the new one: say the Japanese for to wake from sleep, then the Japanese for to lie down.
\end{quote}

\subsection*{You'll want to know: \ja{とこのま}}
\textbf{\ja{とこのま}} — \emph{tokonoma} — \textquotedblleft{}an alcove (for flowers or a hanging scroll)\textquotedblright{}.

\begin{quote}
\textbf{\ja{とこのまに} \ja{はなを} \ja{かざりました。}} — \emph{tokonoma ni hana o kazarimashita} — \textquotedblleft{}I put flowers in the alcove.\textquotedblright{}
\end{quote}

\begin{cousinweb}
\textbf{\ja{とこのま}} (\emph{tokonoma}) is a native Japanese word (wago), inherited rather than borrowed from Chinese, built from \emph{toko}, \textquotedblleft{}raised floor\textquotedblright{}, the particle \emph{no} and \emph{ma}, \textquotedblleft{}space, room\textquotedblright{}; its kanji were attached later as kun'yomi (JMdict).
\end{cousinweb}

\begin{grammarlens}[title={What you've built}]
One more word for getting around.
\end{grammarlens}

\subsection*{Your turn}
\begin{itemize}
\raggedright
  \item \emph{Say it:} \emph{tokonoma}
  \item \emph{Say it:} \emph{tokonoma}, once more
  \item \emph{Say it:} \emph{tokonoma ni hana o kazarimashita} — I put flowers in the alcove.
  \item \emph{Recall:} say the Japanese for to wake from sleep, then the Japanese for to lie down, then say \emph{tokonoma} again
\end{itemize}

\begin{tcolorbox}[breakable,colback=teal!4,colframe=teal!35!black,title={Before you move on}]
What does \ja{とこのま} mean? (\textquotedblleft{}An alcove (for flowers or a hanging scroll)\textquotedblright{}.) Say it once more.
\end{tcolorbox}
\section[\texorpdfstring{\ja{ざしき} (zashiki)}{ざしき (zashiki)}]{\ja{ざしき} (zashiki) — a tatami room}

\label{lesson:JA-C400-zashiki}



\begin{quote}
Before the new one: say the Japanese for to lie down, then the Japanese for an alcove.
\end{quote}

\subsection*{You'll want to know: \ja{ざしき}}
\textbf{\ja{ざしき}} — \emph{zashiki} — \textquotedblleft{}a tatami room\textquotedblright{}.

\begin{quote}
\textbf{\ja{ざしきで} \ja{みんなと} \ja{ごはんを} \ja{たべました。}} — \emph{zashiki de minna to gohan o tabemashita} — \textquotedblleft{}We all had a meal in the tatami room.\textquotedblright{}
\end{quote}

\begin{cousinweb}
\textbf{\ja{ざしき}} (\emph{zashiki}) is built from both layers of the vocabulary: the Sino-Japanese \emph{za}, \textquotedblleft{}seat\textquotedblright{} (Middle Chinese \emph{dzwaH}), an on'yomi borrowed from Chinese, and the native Japanese (wago) \emph{shiki}, \textquotedblleft{}spreading\textquotedblright{} (from \emph{shiku}), inherited rather than borrowed (JMdict; KANJIDIC2; Middle Chinese in Baxter's transcription).
\end{cousinweb}

\begin{grammarlens}[title={What you've built}]
One more word for getting around.
\end{grammarlens}

\subsection*{Your turn}
\begin{itemize}
\raggedright
  \item \emph{Say it:} \emph{zashiki}
  \item \emph{Say it:} \emph{zashiki}, once more
  \item \emph{Say it:} \emph{zashiki de minna to gohan o tabemashita} — We all had a meal in the tatami room.
  \item \emph{Recall:} say the Japanese for to lie down, then the Japanese for an alcove, then say \emph{zashiki} again
\end{itemize}

\begin{tcolorbox}[breakable,colback=teal!4,colframe=teal!35!black,title={Before you move on}]
What does \ja{ざしき} mean? (\textquotedblleft{}A tatami room\textquotedblright{}.) Say it once more.
\end{tcolorbox}
\section[\texorpdfstring{\ja{えんがわ} (engawa)}{えんがわ (engawa)}]{\ja{えんがわ} (engawa) — a veranda (along a Japanese house)}

\label{lesson:JA-C400-engawa}



\begin{quote}
Before the new one: say the Japanese for an alcove, then the Japanese for a tatami room.
\end{quote}

\subsection*{You'll want to know: \ja{えんがわ}}
\textbf{\ja{えんがわ}} — \emph{engawa} — \textquotedblleft{}a veranda (along a Japanese house)\textquotedblright{}.

\begin{quote}
\textbf{\ja{えんがわで} \ja{おちゃを} \ja{のみました。}} — \emph{engawa de ocha o nomimashita} — \textquotedblleft{}We drank tea on the veranda.\textquotedblright{}
\end{quote}

\begin{cousinweb}
\textbf{\ja{えんがわ}} (\emph{engawa}) is built from both layers of the vocabulary: the Sino-Japanese \emph{en}, \textquotedblleft{}edge\textquotedblright{} (Middle Chinese \emph{ywenH}), an on'yomi borrowed from Chinese, and the native Japanese (wago) \emph{gawa}, \textquotedblleft{}side\textquotedblright{}, inherited rather than borrowed (JMdict; KANJIDIC2; Middle Chinese in Baxter's transcription).
\end{cousinweb}

\begin{grammarlens}[title={What you've built}]
One more word for getting around.
\end{grammarlens}

\subsection*{Your turn}
\begin{itemize}
\raggedright
  \item \emph{Say it:} \emph{engawa}
  \item \emph{Say it:} \emph{engawa}, once more
  \item \emph{Say it:} \emph{engawa de ocha o nomimashita} — We drank tea on the veranda.
  \item \emph{Recall:} say the Japanese for an alcove, then the Japanese for a tatami room, then say \emph{engawa} again
\end{itemize}

\begin{tcolorbox}[breakable,colback=teal!4,colframe=teal!35!black,title={Before you move on}]
What does \ja{えんがわ} mean? (\textquotedblleft{}A veranda (along a Japanese house)\textquotedblright{}.) Say it once more.
\end{tcolorbox}
\section[\texorpdfstring{\ja{おてあらい} (otearai)}{おてあらい (otearai)}]{\ja{おてあらい} (otearai) — a toilet; a restroom}

\label{lesson:JA-C400-otearai}



\begin{quote}
Before the new one: say the Japanese for a tatami room, then the Japanese for a veranda.
\end{quote}

\subsection*{You'll want to know: \ja{おてあらい}}
\textbf{\ja{おてあらい}} — \emph{otearai} — \textquotedblleft{}a toilet; a restroom\textquotedblright{}.

\begin{quote}
\textbf{\ja{おてあらいは} \ja{どこですか。}} — \emph{otearai wa doko desu ka} — \textquotedblleft{}Where is the toilet?\textquotedblright{}
\end{quote}

\begin{cousinweb}
\textbf{\ja{おてあらい}} (\emph{otearai}) is a native Japanese word (wago), inherited rather than borrowed from Chinese, built from polite \emph{o-}, \emph{te}, \textquotedblleft{}hand\textquotedblright{}, and \emph{arai}, \textquotedblleft{}washing\textquotedblright{} (from \emph{arau}); its kanji were attached later as kun'yomi (JMdict).
\end{cousinweb}

\begin{grammarlens}[title={What you've built}]
One more word for getting around.
\end{grammarlens}

\subsection*{Your turn}
\begin{itemize}
\raggedright
  \item \emph{Say it:} \emph{otearai}
  \item \emph{Say it:} \emph{otearai}, once more
  \item \emph{Say it:} \emph{otearai wa doko desu ka} — Where is the toilet?
  \item \emph{Recall:} say the Japanese for a tatami room, then the Japanese for a veranda, then say \emph{otearai} again
\end{itemize}

\begin{tcolorbox}[breakable,colback=teal!4,colframe=teal!35!black,title={Before you move on}]
What does \ja{おてあらい} mean? (\textquotedblleft{}A toilet; a restroom\textquotedblright{}.) Say it once more.
\end{tcolorbox}
\section[\texorpdfstring{\ja{しめる} (shimeru)}{しめる (shimeru)}]{\ja{しめる} (shimeru) — to shut (a door or window)}

\label{lesson:JA-C400-shimeru}



\begin{quote}
Before the new one: say the Japanese for a veranda, then the Japanese for a toilet.
\end{quote}

\subsection*{You'll want to know: \ja{しめる}}
\textbf{\ja{しめる}} — \emph{shimeru} — \textquotedblleft{}to shut (a door or window)\textquotedblright{}.

\begin{quote}
\textbf{\ja{さむいので、まどを} \ja{しめて} \ja{ください。}} — \emph{samui node, mado o shimete kudasai} — \textquotedblleft{}It is cold, so please shut the window.\textquotedblright{}
\end{quote}

Polite \textbf{\ja{しめます}}, \emph{shimemasu}; an ichidan verb. Past (polite): \textbf{\ja{ははが} \ja{まどを} \ja{しめました。}}, \emph{haha ga mado o shimemashita}.

\begin{cousinweb}
\textbf{\ja{しめる}} (\emph{shimeru}) is a native Japanese word (wago), inherited rather than borrowed from Chinese; its kanji was attached later as a kun'yomi (JMdict). It is the partner of \emph{shimaru}, \textquotedblleft{}to close\textquotedblright{} (JMdict).
\end{cousinweb}

\begin{grammarlens}[title={What you've built}]
One more word for getting around.
\end{grammarlens}

\subsection*{Your turn}
\begin{itemize}
\raggedright
  \item \emph{Say it:} \emph{shimeru}
  \item \emph{Say it:} \emph{shimeru}, once more
  \item \emph{Say it:} \emph{samui node, mado o shimete kudasai} — It is cold, so please shut the window.
  \item \emph{Recall:} say the Japanese for a veranda, then the Japanese for a toilet, then say \emph{shimeru} again
\end{itemize}

\begin{tcolorbox}[breakable,colback=teal!4,colframe=teal!35!black,title={Before you move on}]
What does \ja{しめる} mean? (\textquotedblleft{}To shut (a door or window)\textquotedblright{}.) Say it once more.
\end{tcolorbox}
